\documentclass[10pt,twocolumn,letterpaper]{article}
\usepackage[rebuttal]{iccv}

% Include other packages here, before hyperref.
\usepackage{graphicx}
\usepackage{amsmath}
\usepackage{amssymb}
\usepackage{booktabs}

% Import additional packages in the preamble file, before hyperref
\input{preamble}

% If you comment hyperref and then uncomment it, you should delete
% egpaper.aux before re-running latex.  (Or just hit 'q' on the first latex
% run, let it finish, and you should be clear).
\definecolor{iccvblue}{rgb}{0.21,0.49,0.74}
\usepackage[pagebackref,breaklinks,colorlinks,allcolors=iccvblue]{hyperref}

% If you wish to avoid re-using figure, table, and equation numbers from
% the main paper, please uncomment the following and change the numbers
% appropriately.
%\setcounter{figure}{2}
%\setcounter{table}{1}
%\setcounter{equation}{2}

% If you wish to avoid re-using reference numbers from the main paper,
% please uncomment the following and change the counter value to the
% number of references you have in the main paper (here, 100).
%\makeatletter
%\apptocmd{\thebibliography}{\global\c@NAT@ctr 100\relax}{}{}
%\makeatother

%%%%%%%%% PAPER ID  - PLEASE UPDATE
\def\paperID{11462} % *** Enter the Paper ID here
\def\confName{ICCV}
\def\confYear{2025}

\begin{document}

%%%%%%%%% TITLE - PLEASE UPDATE
\title{
\vspace{-0.2cm}
Empowering SFDA via MLLM-Guided Reliability-Based Curriculum Learning
\vspace{-0.2cm}
}  % **** Enter the paper title here

\maketitle
\thispagestyle{empty}
\appendix

% \vspace{-2mm}

%%%%%%%%% BODY TEXT - ENTER YOUR RESPONSE BELOW
\noindent 
% We thank all reviewers for their thoughtful feedback. We appreciate the recognition of our paper’s structured design, empirical performance, and technical novelty (Reviewer \textcolor{purple}{wqw3}, \textcolor{orange}{WBNd}), as well as the strong gains enabled by our use of frozen MLLMs (Reviewer \textcolor{blue}{9xwq}). 
 We thank all reviewers for their thoughtful feedback. We appreciate recognition of our framework, contributions (\textcolor{purple}{wqw3}, \textcolor{orange}{WBNd}), and strong empirical performance (\textcolor{blue}{9xwq}).

% \noindent
\textbf{Contribution: }
Our work builds on recent SFDA trends that use frozen vision-language models (e.g., CLIP in PADCLIP [15], DIFO [41]) and extends them with a simple, effective framework that leverages multiple frozen MLLMs for richer supervision. RCL performs multi-teacher distillation without tuning, using agreement-driven, reliability-based partitioning across three stages (RKT, SMKE, MMR). It integrates STS to align open-ended MLLM outputs to class labels and uses a reliability score for pre-adaptation trust estimation. Fig. 6 shows that RCL consistently outperforms both weaker MLLMs and their ensemble, and remains robust across settings. Tab. 4 further highlights that even without MLLMs, RCL achieves competitive or stronger results than prior SFDA methods. Despite minimal adaptation cost (Tab. 7), RCL scales well with MLLM strength. Our supplementary (Sec. 10, Tab. 9–10) further analyzes MMR behavior and confidence thresholds, confirming robustness under varied conditions.

% Our work follows recent SFDA trends that use frozen foundation models (e.g., CLIP in PADCLIP [15], DIFO [41]) and extends them with a simple, effective framework for using multiple frozen MLLMs (now increasingly used across domains). MLLMs are a natural progression from CLIP, offering richer supervision. RCL performs multi-teacher distillation without tuning, using agreement-driven, reliability-based partitioning across three stages (RKT, SMKE, MMR). It integrates STS for aligning open-ended MLLM outputs to class labels and uses a reliability score for pre-adaptation trust estimation. Fig. 6 shows RCL improves even with smaller MLLMs, demonstrating scalability, robustness, and strong performance across settings, all with minimal adaptation cost (Tab. 7). We address reviewer concerns below. Our supplementary (Sec. 10, Tab. 9–10) further analyzes confidence thresholds and MMR behavior, confirming robustness under varied settings.



\noindent \textbf{\textcolor{purple}{\underline{Reviewer wqw3: }}}

\textcolor{purple}{\textbf{Q1:}} \textit{Strengths not emphasized in intro.} 
We agree and will revise the introduction to better foreground our structured, reliability-based curriculum using frozen MLLMs.
% Thank you—we agree and will revise the introduction to foreground our key contributions: (1) a reliability-based curriculum using frozen, agreement-driven MLLMs, and (2) a three-phase training scheme that adapts supervision and MLLM guidance based on confidence. 
% Unlike prior works, our method enables structured multi-teacher adaptation from MLLMs.

\textcolor{purple}{\textbf{Q2:}} \textit{Thresholding and masking may limit novelty.}
While these methods are simple, we integrate them in a structured and progressive curriculum that drives strong results. This mirrors trends in recent work (e.g., DIFO~[41]) which use simple components paired with foundation models.


\textcolor{purple}{\textbf{Q3:}} \textit{Lack of pseudo-label accuracy analysis.} 
Pseudo-label trends in Fig. 10 and Supp. Lines 911–925 validate Eq. (1), (6), and (8). We will add similar plots for weaker MLLMs (Tab. 7) and clarify in the main text.
% Thank you—we agree this is important. We show pseudo-label quality trends in Fig. 10 and provide supporting analysis in Lines 911–925 of the supplementary. These validate our design in Eq. (1), (6), and (8). In Table 7, Setting 1, where weaker MLLMs result in more disagreement, RCL still improves performance. We will add similar pseudo-label accuracy plots for this setting and clarify this linkage in the main paper.

\textcolor{purple}{\textbf{Q4:}} \textit{Model-specific bias from MLLMs.} We address this in Tab. 11, which ablates MLLM combinations and observe stable performance. Fig. 7 shows RCL outperforms individual models and their ensemble, suggesting that gains come from the framework design, not any specific MLLM. We will clarify and acknowledge remaining limitations.


% \textcolor{purple}{\textbf{Q5:}} \textit{Minor issues.} Thanks, we will fix typos and reorder figures accordingly.



\noindent \textbf{\textcolor{orange}{\underline{Reviewer WBNd}}}

\textcolor{orange}{\textbf{Q1:}} \textit{Relies more on MLLMs, SFDA alignment.}
We clarify that our method follows the SFDA protocol: we initialize from a source-pretrained model, then adapt solely using unlabeled target data with frozen MLLMs for guidance—mirroring SOTA methods (similar setup to DIFO).
% We clarify that our method follows the standard SFDA protocol: the source model is used for initialization, and adaptation uses only unlabeled target data. This is consistent with recent SOTA SFDA works like DIFO, which also distill from frozen external models.
% Our approach extends this trend with a structured, tuning-free curriculum over multiple MLLMs, while remaining fully source-free.

\textcolor{orange}{\textbf{Q2:}} \textit{Are performance gains mostly due to large MLLMs?}
We address this empirically. RCL shows strong performance even with weaker MLLMs and their ensemble (Fig. 6), and without MLLMs (Tab. 4), rivaling SOTA baselines which confirms that gains stem from our framework.

\textcolor{orange}{\textbf{Q3:}} \textit{Is evaluation on SFDA benchmarks appropriate?}
Yes, SFDA requires no access to source data during adaptation, which we follow. Using frozen external models for guidance is now common in SFDA, as seen in recent works like PADCLIP, DIFO, evaluated on the same benchmarks.

\textcolor{orange}{\textbf{Q4:}} \textit{+9.4\% improvement misleading.} 
The +9.4\% refers to gains over current SOTA methods w/o MLLMs. RCL also outperforms individual MLLMs (e.g., \textbf{+3.0\%} over best MLLM LLaVA-34B, Tab. 1) and their ensemble (\textbf{+4.3\% }and \textbf{+1.6\%} over ensemble setting 1 \& 2 resp.; Fig. 6), showing that improvements stem from our framework, not ensembling.
We will revise the text accordingly.

% RCL also beats MLLM ensembles (e.g., +6.0% over BLIP2, +4.8% over LLaVA-7B, Tab. 4), confirming gains stem from our framework, not ensembling.
% \textcolor{orange}{\textbf{Q5:}} \textit{Is a comparison with simple MLLM ensemble missing?} This comparison is provided in both Table 4 and Figure 6. RCL clearly outperforms the ensemble average, confirming that the improvement comes from our curriculum design, not from ensembling alone. We will emphasize this more directly in the revision.
 
\noindent \textbf{\textcolor{blue}{\underline{Reviewer 9xwq}}}

\textcolor{blue}{\textbf{Q1:}} \textit{Similarity to DAPL and C-SFDA.}
While we build on established components, RCL differs in key ways. RKT uses inter-MLLM semantic agreement (vs. intra-model confidence in DAPL) to drive curriculum learning. SMKE blends MLLM and student predictions via soft confidence (vs. hard thresholds in C-SFDA). MMR introduces multi-hot masking that utilizes agreement signals across multiple teacher MLLMs (vs. top-k logits or single-model preds). Our framework avoids MLLM prompt tuning or finetuning, adapts purely via reliability structure. These design choices yield clear gains: RCL outperforms DAPL (+6.4\%) and C-SFDA (+16.7\%) quantitatively (Tab. 1 \& 2). We believe this is the first work to unify multiple frozen MLLMs under a structured knowledge adaptation framework. 



\iffalse
\textcolor{blue}{\textbf{Q1:}} \textit{Similarity to DAPL and C-SFDA.}
Thank you for the question. Recent SFDA methods like DAPL, PADCLIP, and DIFO (CVPR 2024) combine established tools (e.g., prompt tuning, contrastive loss) in new frameworks. We follow this trend but differ in key ways. RKT introduces a reliability score via frozen MLLM agreement, guiding a curriculum with no tuning or prompts. Unlike DAPL, our reliability score comes from inter-MLLM semantic agreement rather than intra-model confidence. Unlike C-SFDA, SMKE mixes teacher and student signals based on soft confidence, not hard thresholds. MMR uniquely uses multi-hot masking across MLLMs, not top-k logits. SMKE combines soft MLLM predictions with target confidence for mid-reliability regions, unlike the hard thresholds used in C-SFDA. We believe MLLMs are a natural progression from CLIP, and RCL shows how to effectively and scalably use them in SFDA. 
\fi


\textcolor{blue}{\textbf{Q2:}} \textit{MMR only augmentation.}
MMR goes beyond standard augmentation. It introduces multi-hot class masking and soft pseudo-label alignment to handle low-confidence unlabeled samples. Strong/weak views are used to enforce consistency loss, enabling self-correction. While these ideas are standard in unsupervised learning, our novelty lies in how they are combined in RCL to recover utility from difficult samples in a structured SFDA framework.

\textcolor{blue}{\textbf{Q3:}} \textit{Instruction following failures }
% Instruction-following failures 
refer to MLLM outputs like “Audi” instead of “car,” violating class constraints. STS resolves this by semantically aligning free-form MLLM text to fixed class labels. 
% For pre-adaptation performance estimation, we propose the reliability score, which offers a low-cost proxy for label quality and supports sample-level curriculum structuring. 

\textcolor{blue}{\textbf{Q4:}} \textit{A100 usage concern.}
We used an A100 solely due to availability; it is not required. Only pseudo-label caching from the largest MLLM may benefit from it, smaller MLLMs run on standard GPUs, and adaptation remains lightweight (RN50, Tab. 7). RCL still performs well with smaller MLLMs (Tab. 6) and even with a lightweight RN18 backbone (Tab. 5, line 530).

% As noted in Contribution, MLLM inference is one-time cached and training is lightweight (ResNet-50, Tab. 7). RCL performs well even with smaller MLLMs not needing A100s.

% Thanks for the feedback. The A100 is used only once to run batched MLLM inference, which we cache before training. All adaptation uses a lightweight ResNet-50 (~5ms/sample, Table 7). As shown in Fig. 6, RCL performs well even with smaller MLLMs (Setting 1), showing that our framework is scalable and efficient during adaptation.

\textcolor{blue}{\textbf{Q5:}} \textit{Limited innovation, combines known methods.}
We respectfully disagree. Components like agreement and augmentation exist, our novelty lies in their structured integration: a three-stage curriculum over frozen MLLMs guided by reliability. RCL outperforms current SOTA (Tab.1,2), individual MLLMs and their ensemble (Fig. 6), even with smaller models, showing value beyond simple aggregation.

% We respectfully disagree. While components like agreement and augmentation exist, our contribution lies in their structured integration: a three-stage curriculum over frozen MLLMs guided by sample reliability. RCL consistently outperforms individual MLLMs and their ensemble (Fig. 6), especially with smaller MLLMs—demonstrating real adaptation value beyond simple aggregation.


\iffalse
\textcolor{blue}{\textbf{Q5:}} \textit{Limited innovation from combining known methods.}
We respectfully disagree. While components like agreement and augmentation exist, our contribution lies in their structured integration: a three-stage curriculum over frozen MLLMs that adapts based on sample reliability. RCL consistently improves over both individual MLLMs and their ensemble (Fig. 6), especially in smaller MLLM setting, showing the framework provides real adaptation value, not just aggregation.
\fi


%%%%%%%%% REFERENCES
% {
%     \small
%     \bibliographystyle{ieeenat_fullname}
%     \bibliography{main}
% }

\end{document}
